\documentclass[11pt]{article}
\usepackage[a4paper,margin=28mm]{geometry}
\usepackage{amsmath,amssymb,amsthm,mathtools}
\usepackage[colorlinks=true,linkcolor=blue,citecolor=blue,urlcolor=blue]{hyperref}
\newtheorem{proposition}{Proposition}[section]
\newtheorem{conjecture}[proposition]{Conjecture}
\theoremstyle{definition}
\newtheorem{definition}[proposition]{Definition}
\newtheorem{example}[proposition]{Example}
\newtheorem{remark}[proposition]{Remark}
\newcommand{\ii}{\mathrm i}
\newcommand{\A}{\mathbb A}
\newcommand{\C}{\mathbb C}
\newcommand{\length}{\operatorname{length}}
\newcommand{\Sym}{\operatorname{Sym}}
\newcommand{\ord}{\operatorname{ord}}
\title{Completeness of the Bethe Ansatz From Intersection Theory Aspect}
\author{}
\date{}
\begin{document}
\maketitle

\begin{abstract}
We develop an intersection-theoretic framework for counting solutions of
the Bethe equations for the periodic spin-$1/2$ XXX chain of length $L$
with $n$ Bethe roots, where $L\ge2n$. Completeness in the highest-weight
sector requires the number of physical solutions, counted up to
permutation and with singular solutions appropriately regularized, to
equal $\binom Ln-\binom L{n-1}$. These physical solutions give rise to
Bethe vectors. To count them, we generalize the inclusion--exclusion
principle by constructing a family of auxiliary polynomial systems and
assigning to each ordered solution of the Bethe equations a local index,
defined as a weighted sum of the local lengths of the auxiliary ideals
at that solution. We conjecture that this index equals $1/n!$ at every
physical solution and vanishes at every unphysical solution. Assuming
that solutions of the polynomial $TQ$ relation have no repeated roots
away from $\pm\ii/2$, we establish the conjectured index values for
configurations with collisions among nonsingular roots, configurations
with different numbers of roots at $\ii/2$ and $-\ii/2$, and configurations
containing exactly one singular pair. For configurations with multiple
copies of the singular pair, we express the index as a
normalized dimension of the degree-zero homology of a local complex,
thereby proving its nonnegativity.
\end{abstract}

\section{Background: the XXX Bethe ansatz}

Let $L$ be the length of the periodic spin-$1/2$ XXX chain, with state space
$\mathcal H=(\C^2)^{\otimes L}$. One normalization of its Hamiltonian is
\[
H=\frac14\sum_{\ell=1}^L
\bigl(\boldsymbol{\sigma}_{\ell}\cdot
\boldsymbol{\sigma}_{\ell+1}-1\bigr),
\qquad \boldsymbol{\sigma}_{L+1}=\boldsymbol{\sigma}_1.
\]
The sector with $n$ down spins has dimension $\binom Ln$.
The algebraic Bethe ansatz constructs candidate eigenvectors
\[
B(\lambda_1)\cdots B(\lambda_n)|0\rangle
\]
from parameters satisfying
\begin{equation}\label{eq:bethe}
	\left(\frac{\lambda_k+\ii/2}{\lambda_k-\ii/2}\right)^L
	=
	\prod_{j\ne k}\frac{\lambda_k-\lambda_j+\ii}{\lambda_k-\lambda_j-\ii},
	\qquad k=1,\ldots,n.
\end{equation}
Throughout these notes we assume $L\ge2n$, where $n$ is the number of
Bethe roots. The number of highest-weight states in this
sector, counted with their vector-space dimensions, is
\begin{equation}\label{eq:target}
	d(L,n)=\binom Ln-\binom L{n-1},
\end{equation}
%Indeed, each irreducible $\mathfrak{sl}_2$ summand contributes one state to each of its weight spaces; taking the difference of two consecutive sector dimensions extracts the number of highest-weight states of the specified weight. This is a target dimension for completeness, rather than a count of distinct energy eigenvalues.

One would expect that the Bethe ansatz reconstructs all highest-weight
vectors. For simplicity, rescale the parameters by
\[
\lambda_j=\frac{\ii}{2}x_j,
\qquad x_j=\frac{2\lambda_j}{\ii}.
\]
The singular values $\lambda_j=\pm\ii/2$ become $x_j=\pm1$, and the
interaction shifts $\pm\ii$ become $\pm2$.
After clearing denominators and discarding a nonzero common scalar,
the polynomial is
\begin{equation}\label{eq:polynomial}
	P(x_1,\ldots,x_n)=
	(x_1+1)^L\prod_{j\ne1}(x_1-x_j-2)
	-
	(x_1-1)^L\prod_{j\ne1}(x_1-x_j+2).
\end{equation}
It is symmetric in its last $n-1$ arguments. Set
\[
P_i(x)=P(x_i,x_1,\ldots,\widehat{x_i},\ldots,x_n).
\]
The periodic Bethe equations are $P_i=0$ for $i=1,\ldots,n$.

First we require distinct coordinates; repeated coordinates are called
\emph{collisions}. The Bethe creation operators commute, so we also
identify configurations differing by a permutation. Solutions containing
$x_1=1$ and $x_2=-1$ are called \emph{singular}. For distinct remaining
roots, physical singular solutions of the periodic system require the
additional regularization condition
\begin{equation}\label{eq:physical}
	\left[-\prod_{j=3}^n
	\frac{x_j+1}{x_j-1}\right]^L=1.
\end{equation}
As an example for $n=2$, the empty product is $1$,
and the condition requires $L$ to be even.
See \cite{NW2013} for regularizations and
\cite{HNS2013} for the completeness counting problem.

\section{The counting strategy}
Assume that the set of ordered solutions of $P_1=\cdots=P_n=0$ is finite,
and write
\[
X=\{x^1,\ldots,x^s\}\subseteq\C^n
\]
for this set. In the target count, physical solutions have pairwise
distinct coordinates and yield nonzero Bethe vectors, using
regularization at singular solutions. Each physical permutation orbit
therefore has $n!$ elements. Assign the target weights
\[
w_i=\begin{cases}
	1/n!,&x^i\text{ is a physical solution},\\
	0,&\text{otherwise}.
\end{cases}
\]
The desired completeness statement is
\begin{equation}\label{eq:target-weights}
	\sum_{i=1}^s w_i=d(L,n).
\end{equation}

We work with formal sums of points with rational weights:
\[
Z=\sum_i w_i(Z)[x^i],\qquad w_i(Z)\in\mathbb Q.
\]
Here $[x^i]$ is a formal symbol for the point $x^i$, and equality of
such sums means equality of the weight at every point. We write
$\deg Z=\sum_i w_i(Z)$ for the total weight, allowing negative weights.
All operations below are operations on these formal sums.
For a map $\phi$ of point sets, we use the notation
\[
\phi_*\!\left(\sum_p w_p[p]\right)=\sum_p w_p[\phi(p)],
\]
adding weights when several points have the same image. In
particular, passing from ordered tuples to unordered configurations
simply adds the weights over each permutation orbit. We write
$\Sym^r(\A^1)$ for the set of unordered configurations of $r$ complex
numbers, allowing repetitions.

For a zero-dimensional ideal $I\subseteq\C[x_1,\ldots,x_n]$, define
its global and local lengths by
\[
\begin{aligned}
	\ell(I)&=\dim_\C\bigl(\C[x_1,\ldots,x_n]/I\bigr),\\
	\ell_p(I)&=\dim_\C\bigl(\C[x_1,\ldots,x_n]/I\bigr)_{\mathfrak m_p},
\end{aligned}
\]
where $\mathfrak m_p$ is the maximal ideal of the point $p$, and
$\ell_p(I)=0$ if $p\notin V(I)$. The associated formal sum is
\[
Z(I)=\sum_{p\in V(I)}\ell_p(I)[p],\qquad
\deg Z(I)=\ell(I).
\]
Thus its weights are nonnegative integers, and its total weight
is the dimension of a finite-dimensional quotient algebra.

Our strategy is to express the target weights as rational linear
combinations of ordinary local intersection lengths. We seek ideals
$I_j$ and weights $w_j^{\mathrm{aux}}\in\mathbb Q$ such that each $I_j$ is
generated by $n$ polynomials, is zero-dimensional, and satisfies
$V(I_j)\subseteq X$. Set $Z_j=Z(I_j)$ and write
\[
Z_j=\sum_{i=1}^s l_{ij}[x^i],\qquad
l_{ij}=\ell_{x^i}(I_j),
\]
with $l_{ij}=0$ when $x^i\notin V(I_j)$. We require the formal sum
\begin{equation}\label{eq:weight-cycle}
	Z_{\text{Bethe}}=\sum_jw_j^{\mathrm{aux}}Z_j
\end{equation}
to equal $\sum_iw_i[x^i]$; equivalently,
\begin{equation}\label{eq:local-weight-decomposition}
	\sum_j w_j^{\mathrm{aux}}l_{ij}=w_i\qquad\text{for every }i.
\end{equation}
Taking degrees then gives
\begin{equation}\label{eq:global-weight-decomposition}
	\sum_iw_i=\sum_jw_j^{\mathrm{aux}}\ell(I_j),\qquad
	\ell(I_j)=\deg Z_j=\sum_i l_{ij}.
\end{equation}

The generators of $I_j$ allow us to study their global
intersection numbers. We defer the compactification in $\mathbb P^n$
and the analysis of contributions at infinity to a later section.

The proposed approach has two steps:
\begin{enumerate}
	\item To remove collision and permutation redundancy we make use of M\"obius inversion on the partition lattice. We call this process {\em depermutation}.
	\item Construct an excess formal sum on the singular-pair locus and subtract it.
	The resulting formal sum should satisfy the pointwise identities
	\eqref{eq:local-weight-decomposition} and hence the degree identity
	\eqref{eq:target-weights}.
\end{enumerate}
The first step supplies explicit formal sums $Z_j$ and rational weights
$w_j^{\mathrm{aux}}$ for depermutation. The second modifies these weights and adds
formal sums on the singular locus so that the resulting local sums are the
target weights $w_i$. The advantage is the passage from the local
identities \eqref{eq:local-weight-decomposition} to the global calculation
\eqref{eq:global-weight-decomposition}.

We use the following notation for the auxiliary formal sums. Write $F_i=P_i$,
and let $A,B\subseteq[n]$ be disjoint, with either $A=B=\varnothing$ or
both $A$ and $B$ nonempty. Let $\pi=\{S_1,\ldots,S_m\}$ be a partition
of $R=[n]\setminus(A\cup B)$, and choose a representative $r_\alpha\in
S_\alpha$ for each block. For $t\in R$, define
the reduced polynomial by
\begin{equation}\label{eq:reduced-fixed-equation}
	F_t^{\mathrm{red}}(x_R)
	:=\frac{\left.F_t(x)\right|_{x_a=1\ (a\in A),\,
			x_b=-1\ (b\in B)}}{(x_t+1)^{|A|}(x_t-1)^{|B|}}.
\end{equation}
This is well-defined: the two terms in the numerator are both divisible
by $(x_t+1)^{|A|}(x_t-1)^{|B|}$, since $L\ge2n$.
For $A=B=\varnothing$, set $F_t^{\mathrm{red}}=F_t$.
The dependence on $A,B$ is suppressed in the notation. Define
\begin{equation}\label{eq:fixed-partition-ideal}
	\begin{aligned}
		I_{A,B;\pi}=\bigl(&x_a-1\ (a\in A),\quad x_b+1\ (b\in B),\\
		&F_{r_\alpha}^{\mathrm{red}}\ (1\le\alpha\le m),\\
		&x_t-x_{r_\alpha}\quad
		(1\le\alpha\le m,\ t\in S_\alpha\setminus\{r_\alpha\})
		\bigr).
	\end{aligned}
\end{equation}
Here one reduced equation $F_{r_\alpha}^{\mathrm{red}}=0$ is retained for each block, together
with the equalities within that block. The ideal has $|A|+|B|$ fixed-root
equations, $m$ Bethe equations, and $n-|A|-|B|-m$ equality equations,
for a total of $n$ specified generators. Symmetry makes it independent
of the chosen representatives. Let
\[
Z_{A,B;\pi}=Z(I_{A,B;\pi})
=\sum_{p\in V(I_{A,B;\pi})}\ell_p(I_{A,B;\pi})[p],
\]
whenever $I_{A,B;\pi}$ is zero-dimensional. Its support is contained in
$X$: the original equations at the fixed roots vanish automatically when
both signs occur, and each remaining original equation is
$(x_t+1)^{|A|}(x_t-1)^{|B|}F_t^{\mathrm{red}}$ after fixing $A,B$. Within each block,
the reduced equations agree after imposing the equalities. We abbreviate
$Z_{\varnothing,\varnothing;\pi}$ to $Z_\pi$.

Our aim is to find rational weights $w_{A,B;\pi}^{\mathrm{aux}}$ such that
\[
Z_{\text{Bethe}}=\sum_{A,B,\pi}w_{A,B;\pi}^{\mathrm{aux}}Z_{A,B;\pi}.
\]
This is the notation $j=(A,B;\pi)$ in
\eqref{eq:local-weight-decomposition}; its local lengths are
$l_{i,(A,B;\pi)}=\ell_{x^i}(I_{A,B;\pi})$.
Throughout the note, $I$ denotes an ideal and $Z$ a formal sum of points.
All auxiliary ideals $I_{A,B;\pi}$ use the reduced equations
\eqref{eq:reduced-fixed-equation}; the associated sums $Z_{A,B;\pi}$
therefore need no additional superscript. The original polynomials
$F_t$ are retained for studying local lengths in the full Bethe system.


\begin{thebibliography}{99}
\bibitem{Stacks}
The Stacks Project Authors,
\emph{The Stacks Project},
\href{https://stacks.math.columbia.edu/tag/062F}{Tag~062F},
\href{https://stacks.math.columbia.edu/tag/035E}{Tag~035E}, and
\href{https://stacks.math.columbia.edu/tag/0DW6}{Section~47.21}.

\bibitem{MTV2011}
E.~Mukhin, V.~Tarasov, and A.~Varchenko,
\emph{Reality property of discrete Wronski map with imaginary step},
Lett.\ Math.\ Phys.\ \textbf{100}, 151--160 (2012),
\href{https://arxiv.org/abs/1102.5368}{arXiv:1102.5368}.

\bibitem{PS1999}
G.~P.~Pronko and Yu.~G.~Stroganov,
\emph{Bethe equations ``on the wrong side of the equator''},
J.\ Phys.\ A: Math.\ Gen.\ \textbf{32}, 2333--2340 (1999),
\href{https://arxiv.org/abs/hep-th/9808153}{arXiv:hep-th/9808153}.

\bibitem{MTV2009}
E.~Mukhin, V.~Tarasov, and A.~Varchenko,
\emph{Bethe algebra of homogeneous XXX Heisenberg model has simple spectrum},
Commun.\ Math.\ Phys.\ \textbf{288}, 1--42 (2009),
\href{https://arxiv.org/abs/0706.0688}{arXiv:0706.0688}.

\bibitem{BK2016}
L.~Bus\'e and A.~Karasoulou,
\emph{Resultant of an equivariant polynomial system with respect to the symmetric group},
J.\ Symbolic Comput.\ \textbf{76}, 142--157 (2016),
\href{https://arxiv.org/abs/1407.2799}{arXiv:1407.2799}.

\bibitem{GJ2020}
E.~Granet and J.~L.~Jacobsen,
\emph{On zero-remainder conditions in the Bethe ansatz},
JHEP \textbf{03}, 178 (2020),
\href{https://arxiv.org/abs/1910.07797}{arXiv:1910.07797}.

\bibitem{HJJ2024}
J.~Hou, Y.~Jiang, and R.-D.~Zhu,
\emph{Spin-$s$ Rational $Q$-system},
SciPost Phys.\ \textbf{16}, 113 (2024),
\href{https://arxiv.org/abs/2303.07640}{arXiv:2303.07640}.

\bibitem{Sturmfels}
B.~Sturmfels,
\emph{Solving Systems of Polynomial Equations},
CBMS Regional Conference Series in Mathematics \textbf{97}, AMS (2002),
\href{https://web.math.berkeley.edu/~bernd/cbms.pdf}{author's manuscript}.

\bibitem{HNS2013}
W.~Hao, R.~I.~Nepomechie, and A.~J.~Sommese,
\emph{Completeness of solutions of Bethe's equations},
Phys.\ Rev.\ E \textbf{88}, 052113 (2013),
\href{https://arxiv.org/abs/1308.4645}{arXiv:1308.4645}.

\bibitem{NW2013}
R.~I.~Nepomechie and C.~Wang,
\emph{Algebraic Bethe ansatz for singular solutions},
J.\ Phys.\ A \textbf{46}, 325002 (2013),
\href{https://arxiv.org/abs/1304.7978}{arXiv:1304.7978}.

\end{thebibliography}
\end{document}
